% Lösungen Einheit 3 von 4 – Gleichungen aufstellen und grafisch oder numerisch lösen (zusammenbau v0.1)
\einheitenkopf{Einheit 3 von 4 · Gleichungen aufstellen und grafisch oder numerisch lösen}

\erg{17}{a) $h(t) = 50$ \quad b) $h(t) = 4$: $t_1 = 2$, $t_2 = 4$ \quad c) Die Lösungen sind die Stellen, an denen sich die Graphen schneiden: $x_1 = -2$, $x_2 = 0$, $x_3 = 2$ \quad d) $s(x) - s(x - 3) = 0{,}8$ \quad e) $P(2 | 2)$ und $Q(4 | 8)$ liegen auf dem Graphen: $x = 2$. Ab Woche 2 kommen in den folgenden zwei Wochen 600 Abonnenten hinzu. \quad f) $r(t) = r(t - 2)$: $t \approx 3{,}16$; Punkte $(1{,}16 | 2{,}60)$ und $(3{,}16 | 2{,}60)$ \quad g) $L(12) = 4$, Hälfte $2{,}0$: $\mathrm{sin}(\ldots) = \frac{1}{3}$, $x \approx 7{,}30$; $0{,}30$ h $\approx 18$ min, also 7:18 Uhr \quad h) $\mathrm{sin}\left(\frac{\pi}{6} t\right) = \frac{1}{2}$: $\frac{\pi}{6} t = \frac{\pi}{6}$ oder $\frac{5\pi}{6}$, plus $2\pi$: $t = 1$, $t = 5$, $t = 13$, $t = 17$ \quad i) $f'(x) = -2x + 6$; mal $x$: $-x^2 + 6x - 4 = -2x^2 + 6x$, $x^2 = 4$, $x_1 = 2$, $x_2 = -2$ \quad j) $P(s | s)$: $f(x) = x$; $Q(-t | t)$: $f(x) = -x$; Seitenlängen $s$ und $t$ als Beträge der Lösungen ($P \approx (2{,}46 | 2{,}46)$, $Q \approx (-1{,}47 | 1{,}47)$); Umfänge $4s$ und $4t$, Verhältnis $4s : 4t = s : t$.}

17f)

\begin{ksys}[xmin=0,xmax=8,ymin=0,ymax=4,klein] \funktion{4*\x*exp(-0.5*\x)}{r} \punkt{1.16}{2.60}{P} \punkt{3.16}{2.60}{Q} \end{ksys}

\erg{18}{a) $f(x) = g(x)$ mit dem Rechner: $(-2{,}74 | 0{,}25)$, $(1{,}77 | 2{,}43)$}
\erg{19}{a) Fehler: nur die festen Zeitpunkte 0 und 4 eingesetzt statt eines beliebigen Zeitpunkts $x$. Richtig: $p(x + 4) = p(x) + 30$ \quad b) Der Sinus ist symmetrisch zu $x = \frac{\pi}{2}$: $\mathrm{sin}(\pi - x) = \mathrm{sin}(x)$, also ist auch $\pi - 0{,}30$ eine Lösung. \quad c) $20 - 2t = 15 - t$: $t = 5$ h, Höhe $= 10$ cm \quad d) $f(x) = 3$: $x_1 = -1$, $x_2 = 3$}
